\section{Methodology}

\subsection{Overview}

Separate GP2PL into domain-level plan-library generation and
query-specific goal realisation, corresponding to the two problems
defined in Section~3.

Obtain reusable behavioural evidence from representative planning
instances and compile it once into a reusable BDI plan library
\(\mathcal{L}_D\).

Reuse \(\mathcal{L}_D\) for later achievement and temporally extended
goal queries without regenerating the domain-level plans.

Generate only query-specific coordination plans \(\mathcal{Q}_q\) for
each new query and execute the query with
\[
\mathcal{L}_D^q
=
\mathcal{L}_D \cup \mathcal{Q}_q.
\]

Separate the method accordingly into constructing the reusable
domain library and generating query-specific behaviour over that
library.


\subsection{From Generalized Behaviour to BDI Plan Candidates}

Use a generalized-planning method to obtain reusable behaviour from
the training instances \(\mathcal{I}_{\mathrm{train}}\).

[inline def] Represent the resulting generalized-planning evidence as
a finite set
\[
\mathcal{E}
=
\{e_1,\ldots,e_n\},
\]
where each evidence rule has the form
\[
e
=
\langle g,C,\beta\rangle.
\]

Here \(g\) is an atomic achievement-goal template, \(C\) describes the
conditions under which the behaviour applies, and \(\beta\) is a
parameterised sequence of domain actions for achieving \(g\).

Normalize the evidence against the domain vocabulary and type
hierarchy while preserving shared variables and declared constants.

Transform each validated evidence rule into a BDI plan candidate
\[
r_e
=
+!g:C\leftarrow\beta.
\]

[inline def] Let $\mathcal{C}_{E}$ denote the BDI plan candidates obtained
directly from generalized-planning evidence.

Treat these plans as candidates rather than as the final BDI library,
because the available evidence may cover only part of the domain goal
interface \(\Gamma_D\).


\subsection{Completing the Reusable Plan Library}

[inline def] Let
\[
\Gamma_E
=
\{g \mid \langle g,C,\beta\rangle\in\mathcal{E}\}
\]
denote the achievement-goal templates directly supported by the
generalized-planning evidence.

Identify required goal templates
\[
g\in\Gamma_D\setminus\Gamma_E
\]
that are not directly supported by the available evidence.

Use the domain action model to construct reusable BDI plan candidates
for these missing achievement goals.

Allow generated plans to invoke subsidiary achievement goals when
intermediate conditions can themselves be achieved through reusable
plans.

Combine evidence-derived and domain-derived plans into a common
candidate space for constructing the domain-level library.

Select a reusable plan library
\[
\mathcal{L}_D
\]
that covers the required domain goal interface \(\Gamma_D\).

Require every subsidiary achievement-goal call in
\(\mathcal{L}_D\) to resolve to a type-compatible plan in the same
library.

Require selected plans to be executable under their stated contexts
and to establish their declared achievement goals when they complete
successfully.

Require recursive plan structures to make progress toward completion.

Prefer a compact library when multiple plan collections satisfy the
same goal-coverage and execution requirements.

Return no library when these requirements cannot be satisfied.


\subsection{From Natural-Language Queries to Formal Goals}

Receive a controlled natural-language request together with its declared
parameters, parameter types, and the domain vocabulary.

Translate the request into a parameterised goal specification
\[
\varphi_q\in\Phi,
\]
representing either an achievement goal or a temporally extended goal.

Require the translated specification to use only predicates, argument
types, and temporal operators supported by the domain and the goal
language defined in Section~2.3.

Require the translated specification to preserve the declared query
parameters and their typing constraints.


At query invocation, apply a type-compatible binding
\[
\theta_q
\]
that maps the declared parameters to objects in the current problem
instance.

Obtain the bound query
\[
\widehat{q}=(\varphi_q,\theta_q)
\]
used by the query-specific goal-realisation stage.

\subsection{Query-Specific Goal Realisation}

Treat achievement and temporally extended goals through the same
query-realisation process.

[inline def] Define the execution specification for a bound query as
\[
\psi_q
=
\begin{cases}
F(\varphi_q\theta_q),
&
\varphi_q\in\Phi_{\mathrm{ach}},\\[2mm]
\varphi_q\theta_q,
&
\varphi_q\in\Phi_{\mathrm{teg}}.
\end{cases}
\]

Construct a deterministic finite automaton (DFA) for the finite traces
satisfying \(\psi_q\),
\[
\mathcal{D}_q
=
\langle
Q_q^{\mathrm{dfa}},
2^{AP_q},
\delta_q,
q_q^0,
F_q
\rangle,
\]
where \(Q_q^{\mathrm{dfa}}\) is the set of DFA states,
\(q_q^0\) is the initial state,
\(\delta_q\) is the transition function, and
\(F_q\) is the set of accepting states.

Use the current DFA state to represent the progress already made toward
satisfying the bound query.

Identify DFA transitions that make progress toward an accepting state.

[inline def] Represent the conditions required for a progress transition
with guard \(\chi\) as
\[
\mathcal{O}_{\chi}
=
\langle P_{\chi}^{+},P_{\chi}^{-}\rangle,
\]
where \(P_{\chi}^{+}\) contains conditions that must be established and
\(P_{\chi}^{-}\) contains conditions that must remain false.

Reuse achievement plans from \(\mathcal{L}_D\) to establish the
positive conditions required by a progress transition.

Treat negative conditions as constraints on execution rather than as
new negative achievement goals.

Consider the possible effects of selected plans when several conditions
must hold for the same transition.

Order plan executions when necessary so that achieving a later
condition does not invalidate conditions established earlier.

Reject a progress transition when its required conditions cannot be
safely realised using the available domain-level plans.

Generate a finite set of query-specific BDI plans
\[
\mathcal{Q}_q
\]
that dispatches reusable plans according to the current DFA state and
the progress transitions available from that state.

Append the query-specific plans to the existing domain library,
\[
\mathcal{L}_D^q
=
\mathcal{L}_D\cup\mathcal{Q}_q,
\]
without modifying or regenerating the reusable plans in
\(\mathcal{L}_D\).

Update the DFA state as primitive domain actions change the observed
execution state.

Complete the query successfully only when the DFA reaches an accepting
state.
